\documentclass[11pt]{article}
\usepackage[margin=1in]{geometry}
\usepackage{amsmath,amssymb,amsthm,mathtools}
\usepackage[hidelinks]{hyperref}
\newtheorem{theorem}{Theorem}
\newtheorem{lemma}[theorem]{Lemma}
\newtheorem{proposition}[theorem]{Proposition}
\newtheorem{remark}[theorem]{Remark}
\newcommand{\D}{\mathcal D}
\newcommand{\E}{\mathbb E}
\newcommand{\R}{\mathbb R}
\newcommand{\norm}[1]{\left\lVert#1\right\rVert}
\newcommand{\one}{\mathbf 1}
\DeclareMathOperator{\id}{id}
\title{Common $L^1$ basin boundaries for a bistable mean-field interval map}
\author{Alec Kriebel\\\small Independent researcher\\
\small\href{https://orcid.org/0009-0001-9320-500X}{ORCID: 0009-0001-9320-500X}}
\date{3 October 2026\quad (version 1.0; unrefereed preprint)}
\begin{document}
\maketitle
\begin{abstract}
We prove the common-basin-boundary conjecture stated by Keller in the 2009
Oberwolfach mini-workshop report on the spectrum of transfer operators. For the prescribed
two-branch M\"obius interval maps with feedback
$G(m)=A\tanh(Bm/A)$, $0<A\le2/5$ and $6<B\le16$, the basin of the constant
density equals the boundary of each noncentral stable basin in the relative
$L^1$ topology of all probability densities. We prove the stronger intermediate
statement that the central basin is the $L^1$ closure of the finite preimages
of the constant density. A terminal cumulative transport is lifted backward
with self-consistent feedback. A probability-space inverse contraction and
uniform distortion bound make this construction valid for arbitrary
integrable densities, including unbounded densities and densities with zeros.
The global three-limit theorem and boundary seed are prior results of
Bardet, Keller and Zweim\"uller.
\end{abstract}

\section{The original question}
Let $I=[-1/2,1/2]$, with Lebesgue measure $dx$, and let
\[
 \D=\{u\in L^1(I):u\ge0\text{ a.e.},\ \int_Iu\,dx=1\},
 \qquad \phi(u)=\int_I xu(x)\,dx.
\]
All closures and boundaries below are relative to $\D$ with its $L^1$ metric.
Fix $0<A\le2/5$ and $6<B\le16$. For $|r|\le2/5$ define
\[
 f_r(x)=\frac{(r+4)x+r+1}{2rx+2},\qquad
 T_r(x)=\begin{cases}f_r(x),&x<-r/4,\\ f_r(x)-1,&x>-r/4.\end{cases}
\]
The value at the cut is immaterial for densities. Write $P_r$ for the density
transfer operator of $T_r$, so $(T_r)_*(u\,dx)=(P_ru)\,dx$, and set
\[
 G(m)=A\tanh(Bm/A),\qquad F(u)=P_{G(\phi(u))}u.
\]
The question in \cite[pp.~2713--2715, especially p.~2715]{OWR} concerns exactly
this parameter range and the full space $\D$.

We use the following established facts from
\cite[Theorem~2, Propositions~3 and~4]{BKZ}; they are also stated in
\cite[Theorem~2 and the discussion following it]{OWR}.
There are three fixed densities $\one,u_+,u_-$, every orbit in $\D$ converges
in $L^1$ to one of them, and the basins of $u_+$ and $u_-$ are open. Moreover,
$\one$ is in both basin boundaries. The last fact follows from the earlier
boundary theorem on the analytic integral-representation class. Put
\[
 W=\{u:F^nu\longrightarrow\one\text{ in }L^1\},\qquad
 B_\pm=\{u:F^nu\longrightarrow u_\pm\text{ in }L^1\}.
\]
The cited theorem supplies $\D=W\sqcup B_+\sqcup B_-$. In particular $W$ is
closed. The earlier analytic-core boundary theorem does not itself give a
boundary theorem on all $\D$: that core is compact and is not $L^1$ dense.

\begin{theorem}\label{thm:main}
For the model and parameters above,
\[
 W=\overline{\bigcup_{n\ge0}F^{-n}(\{\one\})}^{\,\D}
   =\partial_\D B_+=\partial_\D B_-.
\]
No boundedness, bounded variation, or positive lower bound on the densities
is assumed.
\end{theorem}

The new step is the full-density finite-preimage closure. The inverse
expansion mechanism and invariant integral representations originate in
\cite[Sections~3--4]{BKZ}; we extend the inverse construction to an arbitrary
probability space and give a simple exact continuous estimate. Sections~2--5
prove the theorem. No differentiable stable-manifold theorem on $L^1$ is used.

\section{An open map on all densities}
For $h_r(x)=(x+r/4)/(1+rx)$, direct algebra gives
$h_r^{-1}=h_{-r}$ and $T_r=T_0\circ h_r$ away from the cut. Let $Q_r$ be
pushforward of densities by $h_r$. It is a signed-$L^1$ isometry, with inverse
$Q_{-r}$. The map $(r,u)\mapsto Q_ru$ is jointly strongly continuous:
this follows first for continuous functions from the smooth Jacobians, then
for $L^1$ functions by approximation and the isometry.

\begin{lemma}\label{lem:open}
$F:\D\to\D$ is a continuous open surjection. For either stable basin,
\[
 F^{-1}(\partial_\D B_\pm)=\partial_\D B_\pm.
\]
\end{lemma}
\begin{proof}
Write $F=P_0K$ with $K(u)=Q_{G(\phi(u))}u$. For $v\in\D$ put
$M_v(r)=\int_I h_{-r}(x)v(x)\,dx$. This decreases in $r$, since
\[
 \partial_r h_{-r}(x)=\frac{x^2-1/4}{(1-rx)^2}\le0.
\]
Thus $r-G(M_v(r))$ increases by at least the parameter increment and changes
sign between $-A$ and $A$. Its unique zero $\rho(v)$ satisfies
$|\rho(v)-\rho(w)|\le(B/2)\norm{v-w}_1$, since $|h_{-r}|\le1/2$ and
$\norm{G'}_\infty\le B$. Consequently
$K^{-1}(v)=Q_{-\rho(v)}v$ is continuous. Checking the mean and scalar equation
in both directions proves this is the inverse of $K$; hence $K$ is a
homeomorphism.

For a prescribed $w\in\D$ let $v=P_0w$. On fibers with $v(T_0x)>0$ put
$q(x)=w(x)/v(T_0x)$, and put $q(x)=1$ on the other fibers. Since
\[
 (P_0w)(y)=\tfrac12\bigl(w(y/2-1/4)+w(y/2+1/4)\bigr),
\]
both branch values vanish on almost every zero fiber. It follows that
$0\le q\le2$ and $P_0q=\one$, including there. For $z\in\D$ define
$R_w(z)(x)=q(x)z(T_0x)$. Then
\[
 P_0R_w(z)=z,\quad R_w(v)=w,\quad
 \norm{R_w(z)-w}_1=\norm{z-v}_1.
\]
These positive sections show that $P_0$ is open and onto, in the relative
topology. This proves the first assertion for $F$.

For $B=B_+$ or $B_-$, tail convergence gives $F^{-1}B=B$.
If $x\in\partial B$, continuity puts $F(x)$ in $\overline B$, while
$x\notin B$ implies $F(x)\notin B$. Thus $F(x)\in\partial B$.
Conversely, if $F(x)\in\partial B$, every relative neighborhood $U$ of $x$
has open image $F(U)$ containing $F(x)$, hence meets $B$. Therefore $U$
meets $F^{-1}B=B$; since $x\notin B$, it is a boundary point.
\end{proof}

\section{A self-consistent inverse contraction}
Put $J=[-1/2,3/2]$. The inverse of the lifted branch is
\[
 b_r(z)=f_r^{-1}(z)=\frac{2z-r-1}{r+4-2rz},\qquad
 f'_r(x)=\frac{4-r^2}{2(1+rx)^2},\qquad
 \partial_r b_r(z)=-\frac{1-4b_r(z)^2}{4-r^2}.
\]
On any probability space, for a bounded $J$-valued random variable $Z$, solve
\[
 r=G(\E b_r(Z)),\qquad V(Z)=b_r(Z).                 \tag{1}\label{eq:V}
\]
The residual is increasing by at least the parameter increment, because
$b_r$ decreases in $r$, and has opposite signs at $\pm A$. Thus the root
exists uniquely. In particular its scalar derivative is at least~1.

\begin{lemma}\label{lem:contract}
On every real probability-space $L^2$, for bounded $J$-valued $Z,Z'$,
\[
 \norm{V(Z)-V(Z')}_2\le\frac{21}{25}\norm{Z-Z'}_2.
\]
The constant is uniform over the stated $A,B$ range.
\end{lemma}
\begin{proof}
Interpolate $Z_t=(1-t)Z+tZ'$, and set $X_t=V(Z_t)$, with root $r_t$.
For each fixed $A>0$, scalar implicit differentiation and differentiation
under expectation are valid: all derivatives of $b_r$ are bounded on the
compact rectangle $J\times[-2/5,2/5]$, the feedback is smooth, and the residual
derivative is at least~1. Put $q=1-4X_t^2$, $g=G'(\E X_t)$,
$\beta=g/(4-r_t^2)$ and let $D$ multiply by $f'_{r_t}(X_t)$. Differentiating
(1) gives the identity
\[
 \dot X_t=\left(\id-\frac{\beta q\E}{1+\beta\E q}\right)D^{-1}\dot Z_t.
                                                               \tag{2}
\]
This differentiates a bounded path, not a map on an unrestricted Banach-space
neighborhood.

We first show that the rank-one factor in (2) has squared norm at most
$C=1+\beta/4$. Let $m=\E q$, $s=\E q^2$; then $m^2\le s\le m$.
When $s>m^2$, in the orthonormal basis $\one,(q-m)/\sqrt{s-m^2}$ the factor
has matrix
\[
 L=\begin{pmatrix}(1+\beta m)^{-1}&0\\
 -\beta\sqrt{s-m^2}/(1+\beta m)&1\end{pmatrix}.
\]
It is the identity on the orthogonal complement. The lower diagonal entry
of $C\id-L^*L$ is $\beta/4$, and direct expansion gives
\[
 \det(C\id-L^*L)=
 \frac{\beta^2\bigl[(Cm-1/4)^2+C(m-s)\bigr]}{(1+\beta m)^2}\ge0.
\]
For $\beta>0$ these imply positive semidefiniteness. For $\beta=0$ the factor
is the identity; constant $q$ gives a scalar contraction on constants and
identity on their orthogonal complement. This proves the claimed bound.

Set $a=|r_t|\le2/5$. The self-consistent tanh relation yields
\[
 0\le g=B(1-r_t^2/A^2)\le16-100a^2,\qquad
 \norm{D^{-1}}\le\frac{2+a}{2(2-a)}.
\]
Consequently the squared norm of the derivative in (2) is bounded by
\[
 U(a)=\frac{(2+a)^2}{4(2-a)^2}
       \left(1+\frac{16-100a^2}{4(4-a^2)}\right)
     =\frac{(2+a)(4-13a^2)}{2(2-a)^3}<\frac7{10}.
\]
The strict inequality on the entire closed interval follows from
\[
 7(2-a)^3-5(2+a)(4-13a^2)
   =172(a-13/43)^2+12/43+58a^3>0.
\]
Integrate (2) along the path and use $7/10<(21/25)^2$.
\end{proof}

We also need elementary uniform branch sensitivities, valid throughout the
same rectangle:
\begin{align}
 0<\partial_z b_r&\le\tfrac34,&|\partial_r b_r|&\le\tfrac{25}{96},\notag\\
 |\partial_x\log f'_r(x)|&\le1,&
 |\partial_r\log f'_r(x)|&\le\tfrac{35}{24}.          \tag{3}\label{eq:sens}
\end{align}
Indeed $1+rx\ge4/5$, $f'_r\ge4/3$, and
$|\partial_r b_r|=(1-4b_r^2)/(4-r^2)\le25/96$.
The last two logarithmic derivatives are $-2r/(1+rx)$ and
$-2r/(4-r^2)-2x/(1+rx)$, respectively.

\section{Backward correction of a central orbit}
\begin{proposition}\label{prop:closure}
Every $u\in W$ is an $L^1$ limit of densities $v_n$ with $F^nv_n=\one$.
\end{proposition}
\begin{proof}
Work on the original probability space $(I,u(x)dx)$; write $\E$ for its
expectation. Let $X_0(x)=x$, $X_{j+1}=T_{r_j}(X_j)$,
$r_j=G(\E X_j)$ and $u_j=F^ju$, the density of $X_j$.
Let $a_j\in\{0,1\}$ be its branch label, so
$f_{r_j}(X_j)=X_{j+1}+a_j$. All $X_j$ are bounded; $u$ need not be in $L^2$.
Fix $n\ge1$ and set $\epsilon=\norm{u_n-\one}_1$. Define
\[
 H_n(y)=-\tfrac12+\int_{-1/2}^y u_n(t)\,dt,\qquad Y_n=H_n(X_n).
\]
$H_n$ is absolutely continuous, nondecreasing, onto $I$, fixes both endpoints,
and $\norm{H_n-\id}_\infty\le\epsilon$.
The cumulative transform sends the nonatomic law of $X_n$ to uniform law,
even if $H_n$ has flat intervals.

Recursively for $j=n-1,\ldots,0$ define
\[
 Y_j=V(Y_{j+1}+a_j),\qquad \rho_j=G(\E Y_j).
                                                               \tag{4}\label{eq:back}
\]
The branch-event sublaw of $Y_{j+1}$ is dominated by its total law. If the
latter is absolutely continuous, each such sublaw is too, and its smooth
inverse-branch pushforward is absolutely continuous. Starting from uniform
$Y_n$, induction proves that every $Y_j$ has a density. The parameter in
(4) is exactly its feedback, and $T_{\rho_j}(Y_j)=Y_{j+1}$ almost surely:
away from endpoints, $b_{\rho_j}(Y_{j+1}+a_j)$ belongs to the branch $a_j$.
Thus the density $v_n$ of $Y_0$ satisfies $F^nv_n=\one$. No independence of
branch labels and later variables is assumed.

Uniqueness in (1) also gives $X_j=V(X_{j+1}+a_j)$. The same labels cancel
in the input difference, so Lemma~\ref{lem:contract}, with $\kappa=21/25$,
gives
\[
 \norm{Y_j-X_j}_2\le\kappa^{n-j}\epsilon,
 \quad\Delta_j:=|\rho_j-r_j|\le16\kappa^{n-j}\epsilon,
 \quad \sum_{j<n}\Delta_j\le84\epsilon.              \tag{5}\label{eq:delta}
\]
The terminal estimate holds because $H_n$ displaces by at most $\epsilon$;
the feedback estimate uses Cauchy--Schwarz on a probability space.

Define deterministic maps backward on each old branch by
\[
 H_j(x)=b_{\rho_j}\bigl(H_{j+1}(T_{r_j}(x))+a_j(x)\bigr).
                                                               \tag{6}\label{eq:H}
\]
Here $a_j(x)$ is the label for $T_{r_j}$ at a point in $I$, not a new random
label; evaluated at $X_j$ it is the label used in (4).
Since $H_{j+1}$ fixes the endpoints, the two limits at the old cut in (6)
are both $b_{\rho_j}(1/2)=-\rho_j/4$. Each $H_j$ is therefore continuous,
nondecreasing and onto, fixes the endpoints, and $Y_j=H_j(X_j)$.
It is absolutely continuous: on each branch the inner old map is smooth
bi-Lipschitz, the later map is absolutely continuous, and the outer inverse
branch is smooth Lipschitz. The finitely many branches glue continuously.
This argument permits flats and does not assume that arbitrary compositions
of absolutely continuous maps are absolutely continuous.

Put $d_j=\norm{H_j-\id}_\infty$. Equations (3) and (6) give
\[
 d_j\le\tfrac34d_{j+1}+\tfrac{25}{96}\Delta_j,\qquad d_n\le\epsilon.
\]
Iterating and then summing, using (5), yields
\[
 d_0\le\tfrac{183}{8}\epsilon,\qquad
 \sum_{j<n}d_j\le3\epsilon+\tfrac{25}{24}\sum_{j<n}\Delta_j
                  \le\tfrac{181}{2}\epsilon.          \tag{7}\label{eq:d}
\]
Every fixed finite old cylinder map is smooth bi-Lipschitz onto its image,
so it pulls exceptional derivative-null sets back to null sets. The chain
rule applied on these cylinders gives, for Lebesgue-a.e. $x$,
\[
 H_0'(x)=u_n(X_n(x))R_n(x),\qquad
 R_n(x)=\prod_{j<n}
   \frac{f'_{r_j}(X_j(x))}{f'_{\rho_j}(H_j(X_j(x)))}.
\]
This makes sense for merely integrable $u_n$; it divides by no density.
Equations (3), (5), and (7) show
\[
 |\log R_n|\le\sum_{j<n}d_j+\tfrac{35}{24}\sum_{j<n}\Delta_j
                 \le213\epsilon.                     \tag{8}\label{eq:log}
\]

The integral required next is against Lebesgue measure, not against $u\,dx$.
For the external sequence of old parameters let
$c_n=P_{r_{n-1}}\cdots P_{r_0}\one$. Then $c_n\le2$. For completeness, the
source integral-representation identities \cite[Section~4.1]{BKZ} give
\[
 w_y(x)=\frac{1-y^2/4}{(1-xy)^2},\quad |y|\le2/3,\qquad
 P_rw_y=p_r(y)w_{\sigma_r(y)}+(1-p_r(y))w_{\tau_r(y)},
\]
where
\[
 p_r(y)=\tfrac12-\frac{r+y}{4+ry},\quad
 \sigma_r(y)=\frac{2(y+r)}{(r+1)y+r+4},\quad
 \tau_r(y)=\frac{2(y+r)}{(r-1)y-r+4}.
\]
Direct substitution in the two inverse branches proves the identity.
For $|r|\le2/5$, $|y|\le2/3$ the weight is in $[0,1]$, both images remain
in $[-2/3,2/3]$, and $1/2\le w_y\le2$.
Since $\one=w_0$, induction gives $c_n\le2$.
Transfer duality and monotone truncation therefore imply
\[
 \int_I|u_n(X_n(x))-1|\,dx
   =\int_I|u_n(y)-1|c_n(y)\,dy\le2\epsilon.
\]
Together with (8) this proves
\[
 \norm{H_0'-\one}_1\le2\epsilon e^{213\epsilon}+e^{213\epsilon}-1.
                                                               \tag{9}\label{eq:derivative}
\]

Finally, any absolutely continuous nondecreasing onto map $H:I\to I$
satisfies $H_*(H'dx)=dx$, by applying the fundamental theorem of calculus to
antiderivatives of continuous test functions. Flat intervals are allowed.
Write $\norm{\cdot}_{\rm TV}$ for the full variation norm of a finite signed
measure, equal to the $L^1$ norm for its density. Pushforward contracts this
norm. For continuous $g$, comparison with $g(H)H'dx$ gives
\[
 \norm{H_*(g\,dx)-g\,dx}_{\rm TV}
 \le\omega_g(\norm{H-\id}_\infty)+\norm{g}_\infty\norm{H'-\one}_1,
                                                               \tag{10}
\]
where $\omega_g$ is its modulus of continuity. Although $H_*(g\,dx)$ may
have atoms, (10) is a measure inequality and remains valid. The actual
$H_{0*}(u\,dx)=v_n\,dx$ is already absolutely continuous by (4).
For any continuous approximation $g$ to the fixed $u$, (7), (9), and (10)
yield
\[
 \norm{v_n-u}_1\le2\norm{u-g}_1+
 \omega_g(183\epsilon/8)+\norm{g}_\infty
                  (2\epsilon e^{213\epsilon}+e^{213\epsilon}-1).
\]
As $u\in W$, $\epsilon\to0$. First fix $g$ and let $n\to\infty$, then let
its approximation error tend to zero. This proves $v_n\to u$ in $L^1$,
without any attraction-rate hypothesis.
\end{proof}

\section{Completion, scope and verification}
Every finite preimage of $\one$ belongs to $W$, and $W$ is closed; together
with Proposition~\ref{prop:closure} this proves the closure identity.
The cited boundary seed and Lemma~\ref{lem:open} place every such preimage in
both basin boundaries. Boundary closedness then gives
$W\subseteq\partial_\D B_\pm$. Conversely, no point of either disjoint open
stable basin can be in either boundary, and global three-limit convergence
exhausts $\D$. Thus $\partial_\D B_\pm\subseteq W$, completing
Theorem~\ref{thm:main}.

The result concerns the specified feedback and rectangle; it does not resolve
the separate finite-particle limiting-mixture question, extend to $A=0$ or
$B=6$, or assert a nonlinear $L^1$ differentiable stable manifold.
The proof is analytic. The accompanying verification supplement retains the
submitted exact identities and interval certificate, the sharper polynomial
certificate above, rough-density/zero-fiber controls, source identities,
and a dated bounded priority audit. Numerical controls are labeled
corroboration and are not used to prove the infinite-dimensional theorem.
The source audit distinguishes accessible preprint versions from the final
journal version, whose proof text was not obtained. The original OWR statement
postdates the journal publication and explicitly formulates this conjecture.

\paragraph{AI use and review status.}
AI tools were used extensively for proof development, verification,
literature research, manuscript preparation, and adversarial review.
The independent adversarial reviews were performed by AI agents.
This is an unrefereed preprint and has not undergone independent external
human peer review. Responsibility for its contents rests with the author.

\begin{thebibliography}{9}
\bibitem{BKZ}
J.-B. Bardet, G. Keller, and R. Zweim\"uller,
\emph{Stochastically stable globally coupled maps with bistable thermodynamic
limit}, Communications in Mathematical Physics \textbf{292} (2009),
237--270, \href{https://doi.org/10.1007/s00220-009-0854-9}{doi:10.1007/s00220-009-0854-9}.
Accessible proof version: \href{https://arxiv.org/abs/0812.4040v1}{arXiv:0812.4040v1};
also \href{https://mat.univie.ac.at/~zweimueller/MyPub/bkz.pdf}{author-hosted manuscript}.
\bibitem{OWR}
\emph{Mini-Workshop: Spectrum of Transfer Operators---Recent Developments and Applications},
Oberwolfach Reports \textbf{6} (2009), no.~4, report~49,
\href{https://doi.org/10.4171/OWR/2009/49}{doi:10.4171/OWR/2009/49}.
G. Keller, \emph{Self-consistent Perron-Frobenius operators for globally coupled maps},
pp.~2713--2715.
\end{thebibliography}
\end{document}
